\chapter{Конструкторская часть}

\section{Общая структура компилятора}

Компилятор построен как последовательность классических этапов трансляции:

\begin{enumerate}
  \item лексический анализ;
  \item синтаксический анализ;
  \item построение абстрактного синтаксического дерева;
  \item ведение таблицы символов;
  \item генерация ассемблерного кода;
  \item сборка и компоновка результата.
\end{enumerate}

Внутреннее взаимодействие модулей организовано через общий набор структур, объявленных в заголовочных файлах. Основные глобальные объекты содержат текущий токен, входной и выходной файлы, таблицу символов, параметры генерации кода и флаги командной строки.

Состав модулей компилятора приведен в таблице \ref{tbl:modules}.

\begin{longtable}{|p{4cm}|p{12cm}|}
  \caption{Основные модули компилятора}
  \label{tbl:modules} \\
  \hline
  \textbf{Файл} & \textbf{Назначение} \\
  \hline
  \endfirsthead
  \hline
  \textbf{Файл} & \textbf{Назначение} \\
  \hline
  \endhead
  \texttt{main.c} & Обработка параметров командной строки, запуск компиляции, сборки и компоновки. \\
  \hline
  \texttt{scan.c} & Лексический анализ: чтение символов, удаление комментариев, распознавание литералов, идентификаторов, ключевых слов и операторов. \\
  \hline
  \texttt{decl.c} & Разбор глобальных и локальных объявлений, функций, параметров, \texttt{auto} и \texttt{extrn}. \\
  \hline
  \texttt{stmt.c} & Разбор операторов: составных блоков, \texttt{if}, \texttt{while}, \texttt{switch}, \texttt{case}, \texttt{goto}, \texttt{return}, меток и выражений-операторов. \\
  \hline
  \texttt{expr.c} & Разбор выражений с учетом приоритетов, построение узлов AST для бинарных, унарных, тернарных операций, вызовов функций и индексирования. \\
  \hline
  \texttt{tree.c} & Создание и отладочная печать узлов абстрактного синтаксического дерева. \\
  \hline
  \texttt{sym.c} & Таблица символов: поиск и добавление глобальных, локальных, параметрических символов и меток. \\
  \hline
  \texttt{types.c} & Внутреннее представление машинных слов, адресов и преобразований размеров данных при генерации кода. \\
  \hline
  \texttt{gen.c} & Машинно-независимый обход AST и выбор операций генерации кода. \\
  \hline
  \texttt{cg.c} & Машинно-зависимая генерация x86-64 assembly. \\
  \hline
\end{longtable}

\section{Лексический анализ}

Лексический анализатор получает поток символов из входного файла и возвращает токены. При чтении выполняются следующие действия:

\begin{itemize}
  \item пропускаются пробельные символы;
  \item удаляются комментарии вида \texttt{/* ... */};
  \item распознаются односимвольные и многосимвольные операторы;
  \item числовые литералы, начинающиеся с \texttt{0}, интерпретируются как восьмеричные;
  \item строки и символьные литералы обрабатывают escape-последовательности языка B;
  \item идентификаторы проверяются на совпадение с ключевыми словами.
\end{itemize}

Для реализации однолексемного просмотра вперед используется механизм возврата токена. Он необходим, например, при разборе оператора, начинающегося с идентификатора: это может быть либо метка вида \texttt{name:}, либо выражение.

\section{Таблица символов}

Таблица символов хранит сведения обо всех именах, известных компилятору. Для каждого символа сохраняются:

\begin{itemize}
  \item имя;
  \item представление значения;
  \item структурный вид символа;
  \item класс хранения;
  \item размер объекта или метка конца функции;
  \item смещение локальной переменной или количество параметров функции.
\end{itemize}

Структурные виды символов:

\begin{itemize}
  \item \texttt{S\_VARIABLE} --- переменная;
  \item \texttt{S\_FUNCTION} --- функция;
  \item \texttt{S\_ARRAY} --- вектор;
  \item \texttt{S\_LABEL} --- метка;
  \item \texttt{S\_EXTERN} --- внешнее имя, объявленное через \texttt{extrn}.
\end{itemize}

Классы хранения:

\begin{itemize}
  \item \texttt{C\_GLOBAL} --- глобальные объекты;
  \item \texttt{C\_LOCAL} --- локальные автоматические объекты;
  \item \texttt{C\_PARAM} --- параметры функции.
\end{itemize}

После завершения генерации тела функции локальные символы удаляются, а глобальные определения остаются доступными для следующих функций.

\section{Абстрактное синтаксическое дерево}

Абстрактное синтаксическое дерево используется как промежуточное представление программы. Узел AST содержит:

\begin{itemize}
  \item код операции;
  \item представление результата выражения;
  \item признак использования в rvalue-контексте;
  \item ссылки на левое, среднее и правое поддеревья;
  \item целочисленное значение, идентификатор символа или размер.
\end{itemize}

Три дочерних указателя позволяют единообразно представлять бинарные операции, тернарный оператор, условные операторы и служебные цепочки операторов. Например, для \texttt{if} левое поддерево хранит условие, среднее --- ветку истинности, правое --- ветку \texttt{else}. Для тернарного выражения левое поддерево содержит условие, среднее --- выражение после \texttt{?}, правое --- выражение после \texttt{:}.

\section{Разбор объявлений}

Глобальные объявления разбираются до конца файла. Если после имени следует \texttt{(}, компилятор обрабатывает функцию. Если после имени следует \texttt{[}, список инициализаторов или \texttt{;}, компилятор обрабатывает внешний объект или вектор.

Историческая форма B:

\begin{verbatim}
name;
name ival1, ival2;
name[constant] ival1, ival2;
name(arg1, arg2) { ... }
\end{verbatim}

Локальные объявления \texttt{auto} размещаются в стековом кадре функции. Автоматические векторы резервируют последовательность машинных слов, адресуемую через операции индексирования и разыменования. Объявления \texttt{extrn} добавляют имя в глобальную таблицу как внешнее до его фактического определения или связывания с внешней функцией.

\section{Разбор выражений}

Выражения разбираются методом приоритетного спуска. Для бинарных операторов задан массив приоритетов. Чем выше значение приоритета, тем раньше оператор связывается с операндами. Унарные и постфиксные операции разбираются отдельно.

Поддерживаются следующие классы выражений:

\begin{itemize}
  \item целые, символьные и строковые литералы;
  \item обращение к идентификатору;
  \item взятие адреса и разыменование;
  \item индексирование вектора;
  \item вызов функции по имени или через вычисленное значение;
  \item префиксный и постфиксный инкремент и декремент;
  \item арифметические, битовые, логические и сравнительные операции;
  \item составные присваивания в стиле B;
  \item тернарный условный оператор.
\end{itemize}

Индексирование строится через адресную арифметику. Для выражения \texttt{v[i]} формируется адрес базы, индекс масштабируется по размеру элемента, затем выполняется сложение и разыменование. Благодаря этому поддерживается свойство B, при котором \texttt{v[i]} и \texttt{i[v]} эквивалентны.

\section{Разбор операторов}

Операторы языка строятся в AST и затем обходятся генератором кода. Составной оператор \texttt{\{ ... \}} формирует цепочку \texttt{A\_GLUE}. Оператор \texttt{while} содержит поддерево условия и поддерево тела. Оператор \texttt{switch} хранит выражение выбора и список \texttt{case}-узлов.

Для меток и \texttt{goto} используется таблица символов. При первом использовании метка может быть добавлена как еще не определенная. После завершения разбора функции выполняется проверка, что все метки, на которые были переходы, действительно определены.

\section{Генерация x86-64 кода}

Генерация кода разделена на два уровня. Модуль \texttt{gen.c} обходит AST и вызывает операции загрузки значений, арифметики, переходов, вызовов функций и сохранения результата. Модуль \texttt{cg.c} реализует эти операции через конкретные инструкции x86-64.

Основные решения генератора:

\begin{itemize}
  \item используется синтаксис GNU assembler;
  \item глобальные объекты размещаются в секции данных;
  \item функции помечаются директивой \texttt{.globl};
  \item локальные переменные размещаются относительно \texttt{\%rbp};
  \item временные значения хранятся в ограниченном наборе регистров;
  \item вызовы функций используют соглашение System V AMD64 ABI;
  \item первые аргументы передаются через регистры \texttt{\%rdi}, \texttt{\%rsi}, \texttt{\%rdx}, \texttt{\%rcx}, \texttt{\%r8}, \texttt{\%r9};
  \item результат функции возвращается через \texttt{\%rax}.
\end{itemize}

Для оператора \texttt{switch} генерируется вспомогательная таблица переходов. Значение выражения сравнивается с константами \texttt{case}, после чего управление передается на соответствующую метку или за пределы оператора выбора.

\section{Вывод}

В конструкторской части описана структура компилятора, назначение его модулей, организация таблицы символов, представление AST, разбор объявлений, выражений и операторов, а также принципы генерации x86-64 кода. Выбранная архитектура отделяет машинно-зависимые операции от общего обхода AST, что упрощает сопровождение и проверку транслятора.
